% Sample Applications, Chapter 1 -- Share of Quarterly Results (annex Sec
% 8.3: scenario, inputs, per-block A-heads, Item / Address / Meaning
% tables).  CP-15: the quarterly results worksheet and its pie chart.  The
% shares and their total are typed and replayed through the shipped
% application (examples/quarterly.i23; GC-11: the total is 116%, asserted
% by test-pubs-i123-examples); the pie chart needs /Graph, which this
% version refuses (era_shortfalls I123-P3-02): DRAFT.  Figure
% screens/sa1.k123.
% GATES: test-i123-formula test-i123-range-menu test-pubs-i123-examples test-pubs-screens
\Gates{test-i123-formula, test-i123-range-menu, test-pubs-i123-examples,
  test-pubs-screens}
\chapter{Share of Quarterly Results}

At the end of each quarter, the planning office reports each division's
share of the quarter's results.  This worksheet lists the five divisions,
their shares as percentages, and the total, and prepares the figures for a
pie chart.

\section{What You Enter}
\begin{itemize}
\item The name of each division
\item Each division's share of the quarter's results, as a percentage
\end{itemize}

\section{Entering the Data}
\ix{percentages}
Type the title in A1, the heading \type{Division} in A3 and \type{"Share}
in B3, the division names in A4..A8 and \type{Total} in A9.  Type each share
with a percent sign, as in \type{40\%}: Initech 123 stores it as 0.4.  In B9
enter \type{@sum(b4..b8)}.  Then format B4..B9 as percentages with no
decimal places (\slashcmd{/Range Format Percent}) and widen column A to 14
characters (\slashcmd{/Worksheet Column Set-Width}).
\FlairShot[fig:sa1]{sa1}{The share of quarterly results}

\begin{LotTable}{p{30mm}p{24mm}p{74mm}}
\LotHead{\textbf{Item} & \textbf{Address} & \textbf{Meaning}}
Division & A4..A8 & The five divisions\\
Share & B4..B8 & Each division's share of the quarter's results\\
Total & B9 & The sum of the shares, \texttt{@SUM(B4..B8)}\LotEnd
\end{LotTable}

\section{Interpreting the Results}
The shares and their total, as the cells show them:
\LotExample{examples/quarterly}

\section{Graphing the Results}
\draft{I123-P3-02: /Graph refuses in this version}
The shares are laid out for a pie chart: the division names in A4..A8 as
the X range and the shares in B4..B8 as the A range.  Select
\slashcmd{/Graph Type Pie}, set the X range with \slashcmd{/Graph X} and the
A range with \slashcmd{/Graph A}, and view the chart with
\slashcmd{/Graph View}.

\section{Changing the Template}
To use the worksheet for another quarter, type the new shares over B4..B8;
the total in B9 is recalculated at once.  To add a division, enter it in
row 9 and move the total down a row with \slashcmd{/Move}, then change the
range in the formula with \fkey{Edit}{F2}.
